\documentclass[12pt]{article}
\usepackage[utf8]{inputenc}
\usepackage[T1]{fontenc}
\usepackage{hyperref}
\title{Piru Language Resources for K-12 Curriculum Development}
\author{Compiled by Claude Sonnet 5.5 (claude-sonnet-5-5) for Andr\'as Kornai}
\date{}
\begin{document}
\maketitle

\section{Introduction}
Piru (retired ISO 639-3 code ppr; Glottocode piru1241) was the northern dialect of Luhu, an Austronesian language of western Seram, Maluku, Indonesia. Sources in the search results describe it as probably extinct, with Luhu itself reported extinct by 2010. The code has been retired. A school curriculum is therefore not realistic; only descriptive documentation was found.

\section{Primary Educational Resources}
None found.

\section{Linguistic and Reference Materials}
\begin{itemize}
\item Wikipedia on the Luhu language (Piru as a dialect): \url{https://en.wikipedia.org/wiki/Luhu_language}
\item Wikipedia on the Piru Bay languages: \url{https://en.wikipedia.org/wiki/Piru_Bay_languages}
\item Max Planck numerals database page for Luhu: \url{https://lingweb.eva.mpg.de/channumerals/Luhu.htm}
\item Glottolog reference record, Stokhof 1982: \url{https://glottolog.org/resource/reference/id/472545}
\end{itemize}

\section{Religious and Community Resources}
None found.

\section{Digital and Audio Resources}
\begin{itemize}
\item Endangered Languages Project record (status information only): \url{https://endangeredlanguages.com/lang/2773}
\end{itemize}

\section{Curriculum Development Context}
Because the variety is extinct or nearly so, any educational use would be heritage or reclamation work based on historical documentation; no existing teaching materials are known.

\end{document}